\documentclass[10pt]{article}
\usepackage[margin=1in]{geometry}
\usepackage[T1]{fontenc}
\usepackage{lmodern,microtype,amsmath,amssymb,amsthm}
\theoremstyle{definition}
\newtheorem{assumption}{Assumption}
\begin{document}
\section*{Stationarity with Exact Hessian Access}
We seek approximate stationarity of a possibly nonconvex function $f:\mathbb{R}^d\to\mathbb{R}$, where $d$ is a nonnegative integer. The function is three times continuously differentiable, with third-derivative norm bounded by $L_0+L_1\|\nabla f(x)\|$ at each point $x$, where $L_0,L_1\geq0$ and $L_0+L_1>0$. At the reference point $x^0=0$, the objective gap above the infimum is at most $\Delta>0$ and the gradient norm is at most $G\geq0$. The oracle returns the exact value, gradient, and Hessian. For $\varepsilon>0$, the goal is a measurable output $\bar{x}$ with $\mathbb{E}_{P}\|\nabla f(\bar{x})\|\leq\varepsilon$, where $P$ is the law of the algorithm's randomness.

The underlying optimization problem is
\[
\inf_{x\in\mathbb{R}^d} f(x).
\]
We use the Euclidean norm on $\mathbb{R}^d$. Derivative norms $\|\cdot\|_{\mathrm{op}}$ are the induced operator norms; in particular, $\nabla^2 f(x)$ is a linear operator on $\mathbb{R}^d$, and $\nabla^3 f(x)$ is a trilinear form. The algorithm receives $d,L_0,L_1,\Delta,G$, and $\varepsilon$.

The parameter regime is restricted by
\[
\begin{aligned}
16\varepsilon&\leq G,\\
\rho:=\frac{\Delta\sqrt{L_0+L_1\varepsilon}}
{\varepsilon^{3/2}}&\geq 2^{39},\\
d&\geq 2^{140}(1+\lceil\rho\rceil)^4.
\end{aligned}
\]
Here $\lceil\rho\rceil$ is the least integer greater than or equal to $\rho$.

Our assumptions:
\begin{assumption}[Third-derivative bound]
The function $f$ belongs to $C^3(\mathbb{R}^d;\mathbb{R})$ and satisfies
\[
\|\nabla^3 f(x)\|_{\mathrm{op}}
\leq L_0+L_1\|\nabla f(x)\|,
\qquad x\in\mathbb{R}^d.
\]
\end{assumption}

\begin{assumption}[Initial objective and gradient control]
At the reference point $x^0=0$,
\[
\begin{aligned}
f(x^0)-f(x)&\leq\Delta
&&\text{for every }x\in\mathbb{R}^d,\\
\|\nabla f(x^0)\|&\leq G.
\end{aligned}
\]
\end{assumption}

Queries may be made at any $x\in\mathbb{R}^d$. Each oracle call returns
\[
\bigl(f(x),\nabla f(x),\nabla^2 f(x)\bigr).
\]
After the allotted oracle calls, the output $\bar{x}\in\mathbb{R}^d$ must be measurable under $P$, its gradient norm must be integrable, and it must satisfy
\[
\mathbb{E}_{P}\bigl[\|\nabla f(\bar{x})\|\bigr]
\leq\varepsilon.
\]
\end{document}
